% Manuscript section, not a compiled or machine-checked proof.
\section{The reverse implication inside Peano arithmetic}
\label{sec:reverse-pa}

The upstream formal theorem is
\[
 \mathrm{PA}\vdash \neg H_{196}\longrightarrow\operatorname{Con}(\mathrm{PA}).
\]
Its separate Lean proof of standard-model nonhalting does not, by itself,
establish the converse implication in PA. The converse has a direct
syntactic proof. The purpose of this section is to specify that proof and
to isolate it from the external hereditarily finite model used upstream.
The additional compiler described below has not been implemented or checked
in Lean in this package.

\subsection{A syntactic interpretation into PA}

Use a conservative definitional extension of ordinary PA containing
exponentiation and binary digits. Write $\beta(s,x)$ for the assertion that
the $x$th binary digit of $s$ is one. Two fixed PA theorems are
\[
 \beta(s,x)\longrightarrow x<s,
 \qquad \neg\beta(x,x).
\]
The extension can be eliminated by primitive recursive proof translation;
its consistency predicate is therefore interchangeable with the ordinary
one for this argument.

Translate a decoded command formula into an arithmetic formula as follows:
\[
 \top^\tau=(0=0),\quad
 (x_a\mathrel{\mathrm M}x_b)^\tau=\beta(x_a,x_b),\quad
 (x_a\mathrel{\mathrm E}x_b)^\tau=(x_a=x_b),
\]
and commute with implication and universal quantification. Bound auxiliary
variables in the definition of $\beta$ are renamed away from the surrounding
binders; choose their labels above the largest label of the finite input.
This is a primitive recursive, capture-avoiding substitution operation.
The translation introduces no additional free variables. It is a map on
syntax, not a satisfaction predicate. The decoder is primitive
recursive: it recurses on strictly smaller nonzero Cantor-coded
subexpressions. Every variable label in the decoded formula of code $c$ is
less than $c$. Consequently $x_c$ is fresh in that formula, a fact PA proves
about the decoder.

The special code $0$ requires an explicit convention. The arithmetic
operation $\mathrm I(0,0)$ also gives $0$, while the decoder reads $0$ as
truth. These two readings are propositionally equivalent. A proof compiler
uses this fixed equivalence in the zero case; it must not assert literal
syntactic equality between the two readings. Code $1$ translates to
$\beta(x_0,x_0)$, whose negation has a fixed PA proof.

There is a related normalization for derived connectives. Coded negation
$\mathrm N(c)=\mathrm I(c,1)$ initially translates to
$\tau(c)\to\beta(x_0,x_0)$, not to literal arithmetic negation.
The fixed PA theorem $\forall x\,\neg\beta(x,x)$ gives a uniform proof
that these are equivalent. The same normalization applies to coded
conjunction and existential quantification, including inside binders.
These fixed normalization proofs are part of the axiom templates below;
the argument does not silently decode code $1$ as a primitive falsity
symbol.

\subsection{Uniform finite comprehension is a proof schema}

For every arithmetic formula $\varphi(x,\vec y)$, PA proves
\begin{equation}\label{eq:finite-comprehension-audit}
 \forall n\;\exists s<2^n\;\forall x<n\,
       (\beta(s,x)\leftrightarrow\varphi(x,\vec y)).
\end{equation}
Here is the proof, with its uniformity made explicit. Induct on $n$ using
the displayed existential formula as the induction formula. At $n=0$ take
$s=0$. At the successor step, use excluded middle for
$\varphi(n,\vec y)$. If it is false, retain $s$; if it is true, replace $s$
by $s+2^n$. The standard binary-digit lemmas show that the lower digits are
unchanged and that the new digit has the required value. The size bound
becomes $s'<2^{n+1}$.

This describes a primitive recursive map from the code of $\varphi$ to a
PA proof of~\eqref{eq:finite-comprehension-audit}. It inserts $\varphi$ into
one full-induction instance and fixed propositional and arithmetic proof
templates. PA proves that this map produces valid proofs by checking the
templates and the induction-axiom recognizer. No assertion of the truth of
arbitrary formulas, and no uniform reflection principle, is used.

Take
\[
 \psi(x)= (x=b\lor\beta(a,x))\land\varphi(x)
 \quad\hbox{and}\quad n=a+b+1.
\]
All possible members of $a$, as well as the distinguished point $b$, are
below $n$. Formula~\eqref{eq:finite-comprehension-audit} yields a number $s$
with
\[
 \forall x\;\bigl(\beta(s,x)\leftrightarrow
       ((x=b\lor\beta(a,x))\land\varphi(x))\bigr).
\]
In particular,
\[
 \begin{split}
 &\forall x\;((\beta(a,x)\to\varphi(x)\to\beta(s,x))
       \land(\beta(s,x)\to\varphi(x))),\\
 &\varphi(b)\to\beta(s,b).
 \end{split}
\]
This is the pointed-separation statement needed by the gate.

For the actual gate, $a,b,c$ are variable labels, and the matrix is the
arithmetic translation of the formula with code $c$, viewed as a formula
in the variable $x_b$. The guard $c\ne a,b$ and the decoder freshness
lemma ensure that the set variable $x_c$ is fresh. One must distinguish
the free distinguished-point occurrence of $x_b$ from its bound occurrence
inside the separation clause; rename the latter to a fresh variable before
inserting the comprehension proof. If $a\ne b$, the preceding construction
applies with $x_a$ as the fixed container parameter. If $a=b$, the membership
antecedent under the quantifier becomes $\beta(x,x)$ and is false. In that
case the simpler choice $s=2^{x_b}$ when $\varphi(x_b)$ holds, and
$s=0$ otherwise, proves both the point clause and
$\forall x(\beta(s,x)\to\varphi(x))$. This separate case avoids treating
the captured occurrence of $x_a$ as a fixed parameter. Equality
and membership vacuity are ordinary quantifier theorems under the same
guard. Classical propositional reasoning combines these facts into the
translation of the exact bundled gate.

All six remaining axiom heads translate to fixed logical or equality
theorem schemas. The restricted propositional head translates to a
classical tautology; its numerical restriction does not affect soundness.
The translation respects modus ponens and generalization. Thus there is a
primitive recursive proof compiler $J$ for which PA proves
\begin{equation}\label{eq:compiler-sound-audit}
 \operatorname{Der}_{196}(d,c)\longrightarrow
 \operatorname{Prf}_{\mathrm{PA}}(J(d),\tau(c)).
\end{equation}
When $c=1$, compose with the fixed proof of
$\neg\beta(x_0,x_0)$ to obtain a PA refutation. This is the precise
syntactic replacement for an appeal to external semantic soundness.

More explicitly, define $J$ by recursion on the annotated proof tree.
At an axiom leaf insert its fixed logical/equality proof template, or
the comprehension template just constructed for a gate. At an MP node
append the translated premise proofs and one MP step, inserting the
fixed zero-code equivalence when necessary. At a generalization node
append the corresponding generalization. The usual finite-proof list
concatenation and reindexing operations are primitive recursive. An
induction on the coded proof tree verifies the output proof predicate.
This defines the mathematical compiler; writing its concrete Gödel-code
implementation and verifying that implementation in Lean remain separate
tasks.

\subsection{Command safety under consistency}

For a command store $(t,w,d,h)$, retain proof codes for the translations of
$t$, $h$, and every entry of the finite stack encoded by $w$. The initial
store has zeros everywhere, so fixed proofs of truth supply the invariant.
Stack underflow produces code $0$ and is harmless. Each command preserves
the invariant by primitive recursive operations on proof codes:
\begin{itemize}
\item An axiom command appends proofs supplied by the head and gate
  compilers and by generalization of the existing proof of $h$.
\item A successful modus-ponens test applies modus ponens to the existing
  proofs; a failed test writes zeros, for which the fixed truth proof is
  used.
\item A shift only moves existing proofs and stack entries. Its halt test
  is $t+h=1$, so one of the proved formula codes would be $1$. Under
  $\operatorname{Con}(\mathrm{PA})$, this is impossible by
  \eqref{eq:compiler-sound-audit} and the fixed refutation of its translation.
\end{itemize}
This is an arithmetical invariant involving proof codes. Induction over a
coded command computation is an instance of full PA induction. It works
for internal, possibly nonstandard command computations as well as for
ordinary finite computations.

\subsection{From commands to the literal tape machine}

The final bridge is the two-level positive-continuation argument already
used by the upstream external simulation. Each safe command boundary has
a finite positive register-program continuation to the next safe command
boundary. Each register instruction has a finite positive tape continuation
to its next frame. The arithmetic versions of the pairing, unpairing,
register and tape-macro contracts are total in PA. Their loop proofs use
decreasing natural-number counters; the finite per-site coverage checks
select the applicable contract.

In particular, the existing arithmetic modules
\path{OriginalModelCoverage} and
\path{EndpointInternalCodedSourcePhysicalStep} supply all-model
instruction contracts and positive coded physical continuations.
\path{InternalCodedSourceDispatchTotal} supplies positive coded command
continuations. Their conclusions are stronger than external standard-model
simulation alone.

Project the auxiliary machine traces to the literal 196-state machine by
relabeling their control vectors. The upstream theorem
\path{StateMapTrace.step_maps} proves the pointwise projection for every
nonpadding control. The same finite table check preserves that property
at each next step. All physical root frames have the fixed nonpadding
control $81$, so induction along any positive frame continuation proves
the required projection, with unchanged length and tape vectors. The
packaged theorem \path{StateMapProjection.segment_to_literal} is
specialized to blank-start halting segments; the positive-segment version
just described is an additional composition obligation, not a claim that
this specialized theorem already proves nonhalting. Apply the principle
below to the projected continuations. An independent backwards simulation
is not needed.

For clarity, the elementary logical principle behind this bridge is the
following. Suppose $I$ is an arithmetical configuration invariant, the
initial configuration satisfies $I$, and every $I$-configuration has a
positive finite execution to another $I$-configuration. A deterministic
machine cannot have a finite halting run from an $I$-configuration. Prove
this inside PA by strong induction on the proposed halting length $n$.
For a continuation of length $k>0$, if $k\le n$, the next invariant
configuration would halt in $n-k<n$ steps. If $k>n$, the continuation would
extend strictly beyond the first halt, contradicting the definition of
the transition relation. The induction predicate may contain $I$;
full PA induction permits that predicate.

Applied to the proof-code invariant, the startup trace and the two levels
of macro simulation, this gives the mathematical implication
\[
 \mathrm{PA}\vdash \operatorname{Con}(\mathrm{PA})\longrightarrow\neg H_{196}.
\]
Together with the upstream arithmetized refutation-to-halting compiler it
gives the desired biconditional. Its exact implementation status must remain
visible: the upstream Lean library checks the latter compiler and external
nonhalting; it does not contain the syntactic interpretation compiler $J$
or the foregoing PA-internal proof-code safety theorem. A Lean certification
of the full biconditional requires adding those modules and the positive
segment and invariant compositions just described. No such certification
is supplied here.

\subsection{What ConciseIPRA contributes}

The native scalar calculi in ConciseIPRA are proved equivalent to iPRA,
whose induction is quantifier-free. They do not by that fact furnish a
refutation calculus for full PA. In particular the uniform finite
comprehension proof above uses the full induction instance containing
the arbitrary translated formula $\varphi$. Replacing the present calculus
by an iPRA calculus without supplying the missing strength would change the
consistency problem. Useful transferable ideas are operational: deriving
equality rules from other rules, choosing antecedent order for shared
instruction trees, and using a literal syntactic feature to select a rule
mode. Those ideas are applicable only after the full-PA proof translations
are re-established.
